\documentclass[10pt]{article}
\usepackage[T1]{fontenc}
\usepackage[utf8]{inputenc}
\usepackage[french]{babel}
\usepackage{amsmath,amssymb,mathrsfs,mathpazo}
\usepackage{geometry}
\geometry{a4paper,margin=2cm}
\usepackage{graphicx,array,tabularx,tabularray,colortbl,multirow,multicol,booktabs,enumitem,pifont,ulem,fancybox,diagbox,cancel}
\usepackage{pstricks,pstricks-add,pst-plot,pst-tree,pst-node,pst-text,pst-func,pst-math,pst-bspline,pst-3dplot,pst-eucl,pst-solides3d}
\let\pscalculate\undefined
\usepackage{pst-bezier}
\usepackage{tikz,pgfplots,tkz-tab}
\usetikzlibrary{arrows,arrows.meta,patterns,calc,positioning,trees,decorations.markings,shapes.geometric,shapes.misc,shapes.arrows,intersections,decorations.pathmorphing,decorations.pathreplacing,backgrounds,shadows,babel}
\pgfplotsset{compat=1.18}
\usepgfplotslibrary{fillbetween}
\usepackage[np]{numprint}
\usepackage{listings,makecell,pas-tableur,twoopt}
\let\barre\undefined
\usepackage[pstricks]{tablvar}
\definecolor{darkgreen}{RGB}{0,100,0}
\providecommand{\red}{\color{red}}
\providecommand{\blue}{\color{blue}}
\providecommand{\green}{\color{green}}
\providecommand{\black}{\color{black}}
\providecommand{\white}{\color{white}}
\providecommand{\gray}{\color{gray}}
\providecommand{\cyan}{\color{cyan}}
\providecommand{\magenta}{\color{magenta}}
\providecommand{\yellow}{\color{yellow}}
\providecommand{\violet}{\color{violet}}
\providecommand{\orange}{\color{orange}}
\providecommand{\euro}{\texteuro}
\providecommand{\EUR}{\texteuro}
\providecommand{\ieme}{\textsuperscript{e}}
\providecommand{\degres}{\ensuremath{^\circ}}
\providecommand{\degree}{\ensuremath{^\circ}}
\providecommand{\vv}[1]{\overrightarrow{#1}}
\providecommand{\Coord}[1]{#1}
\providecommand{\decosix}{}
\providecommand{\starredbullet}{\ensuremath{\bullet}}
\pagestyle{empty}
\setlength{\parindent}{0pt}
\usepackage{adjustbox}\DeclareUnicodeCharacter{2236}{:}\DeclareUnicodeCharacter{2190}{\ensuremath{\leftarrow}}
\providecommand{\R}{\mathbb{R}}
\providecommand{\N}{\mathbb{N}}
\providecommand{\Z}{\mathbb{Z}}
\providecommand{\Q}{\mathbb{Q}}
\providecommand{\D}{\mathbb{D}}
\providecommand{\C}{\mathbb{C}}
\providecommand{\e}{{\mathrm{e}}}
\providecommand{\vect}[1]{\overrightarrow{#1}}
\providecommand{\vectt}[1]{\overrightarrow{\mathrm{#1}}}
\providecommand{\vvt}[1]{\overrightarrow{\mathrm{#1}}}
\providecommand{\np}[2][]{{#2~#1}}
\providecommand{\vv}[1]{{\overrightarrow{#1}}}
\renewcommand{\d}{{\mathrm{d}}}
\providecommand{\degres}{{^\circ}}
\providecommand{\degre}{{^\circ}}
\providecommand{\blue}{\color{blue}}
\providecommand{\red}{\color{red}}
\providecommand{\enskip}{\quad}
\providecommand{\llbracket}{\lbrack\!\lbrack}
\providecommand{\rrbracket}{\rbrack\!\rbrack}
\providecommand{\euro}{{\text{€}}}
\providecommand{\no}{{\text{n°}}}
\providecommand{\interval}[2]{{[#1;#2]}}
\providecommand{\intervFO}[2]{{[#1;#2[}}
\providecommand{\dotfill}{\ldots}

\providecommand{\ds}{\displaystyle}
\providecommand{\textding}[1]{#1}
\providecommand{\Oij}{\(\mathrm{O};\vec{\imath},\vec{\jmath}\)}
\providecommand{\Oijk}{\(\mathrm{O};\vec{\imath},\vec{\jmath},\vec{k}\)}
\providecommand{\Ouv}{\(\mathrm{O};\vec{u},\vec{v}\)}
\providecommand{\intervFF}[2]{\ensuremath{[#1;#2]}}
\renewcommand{\euro}{\ifmmode\text{\texteuro}\else\texteuro\fi}
\DeclareUnicodeCharacter{00B0}{\ifmmode^{\circ}\else\textdegree\fi}
\DeclareUnicodeCharacter{20AC}{\ifmmode\text{\texteuro}\else\texteuro\fi}
% Énoncés : Éducation nationale. Transcription et corrigés : APMEP, Denis Vergès et contributeurs. Découpage et classement : CoopMaths, Sébastien Lozano, Éric Elter et contributeurs. Adaptation de la présentation : Maths Explorer.
% CC BY-SA : https://coopmaths.fr/www/about/
% https://coopmaths.fr/alea/static/e3c/2020/tex/e3c_2020_00_sujet18_4_cor.tex
\begin{document}
\noindent\textit{APMEP / CoopMaths. CC BY-SA. Présentation adaptée par Maths Explorer.}\par\medskip
\subsection*{1.}
	
	On considère la fonction \(f\) définie par :
	\[
	f(x) = x^3 + 3x^2 - 9x - 1.
	\]
	
	\subsubsection*{a.}
	
	On a :
	\[
	f'(x) = 3x^2 + 6x - 9 = 3(x^2 + 2x - 3).
	\]
	
	\subsubsection*{b.}
	
	Le signe de \(f'(x)\) est celui du trinôme \(x^2 + 2x - 3\).
	
	Pour celui-ci :
	\[
	\Delta = 2^2 - 4 \times 1 \times (-3) = 4 + 12 = 16 = 4^2 > 0.
	\]
	
	Le trinôme a donc deux racines :
	\[
	x_1 = \dfrac{-2 + 4}{2} = 1 \quad \text{et} \quad x_2 = \dfrac{-2 - 4}{2} = -3.
	\]
	
	On sait que ce trinôme est positif sauf sur l'intervalle \(\left] -3 ; 1 \right[\) où il est négatif.
	
	Il en résulte que la fonction \(f\) est croissante sauf sur l'intervalle \(\left] -3 ; 1 \right[\) où elle est décroissante.
	
	Avec \(f(-3) = 26\) et \(f(1) = -6\), on a le tableau de variations :
	
	\begin{center}
	\begin{tikzpicture}[double distance=2pt]
		\tkzTabInit{$x$/1,$f'(x)$/1,$f(x)$/3}{$-\infty$,$-3$,$1$, $+\infty$}
		\tkzTabLine{,+,z,-,z, +}
		\tkzTabVar{-/,+/$26$,-/$-6$,+/}
	\end{tikzpicture}
\end{center} 

	
	\subsubsection*{c.}
	
	On sait qu'une équation réduite de la tangente à \(C_f\) au point d'abscisse \(-1\) est :
	\[
	y - f(-1) = f'(-1)(x - (-1)).
	\]
	
	Avec \(f(-1) = 10\) et \(f'(-1) = -12\), l'équation devient :
	\[
	y - 10 = -12(x + 1) \quad \Leftrightarrow \quad y = -12x - 12 + 10 \quad \Leftrightarrow \quad y = -12x - 2.
	\]
	
	\subsection*{2.}
	
	\subsubsection*{a.}
	
	La parabole a sa concavité tournée vers le haut : on a donc \(a > 0\) (on rappelle que \(a \neq 0\)).
	
	La parabole n'a pas de point commun avec l'axe des abscisses, donc le trinôme n'a pas de racines : \(\Delta < 0\).
	
	\subsubsection*{b.}
	
	Un point est commun si et seulement si ses coordonnées vérifient les équations des deux courbes ; il faut donc résoudre l'équation :
	\[
	x^3 + 3x^2 - 9x - 1 = 10x^2 + 8x + 8 \quad \text{ou} \quad x^3 - 7x^2 - 17x - 9 = 0.
	\]
	
	Ceci n'est pas possible, mais sur la figure, on voit que les deux courbes ont un point commun d'abscisse \(-1\) :
	\[
	f(-1) = 10 \quad \text{et} \quad g(-1) = 10 - 8 + 8 = 10.
	\]
	
	De plus, \(f'(-1) = -12\) et comme \(g'(x) = 20x + 8\), on a \(g'(-1) = -20 + 8 = -12\).
	
	Au point d'abscisse \(-1\), les deux fonctions ont le même nombre dérivé, donc la même tangente.

\end{document}
